\chapter{2008年北京卷（理）}

\section{单选题}

\begin{problem}
已知全集 \(U = \R\)，集合 \(A = \{x \mid -2 \leqslant x \leqslant 3\}\)，\(B = \{x \mid x < -1\text{\ 或\ }x > 4\}\)，那么集合 \(A \cap (\complement_U B)\) 等于（\quad）

\choices
    {\(\{x\mid -2\leqslant x < 4\}\)}
    {\(\{x\mid x\leqslant 3\text{\ 或\ }x\geqslant 4\}\)}
    {\(\{x\mid -2\leqslant x < -1\}\)}
    {\(\{x\mid -1\leqslant x\leqslant 3\}\)}

\begin{answer}
D
\end{answer}

\begin{solution}
因为 \(B=(-\infty,-1)\cup(4,+\infty)\)，所以 \(\complement_U B=[-1,4]\). 因而 \(A\cap(\complement_U B)=[-2,3]\cap[-1,4]=[-1,3]\)，对应选项 D.
\end{solution}
\end{problem}

\begin{problem}
若 \(a = 2^{0.5}\)，\(b = \log_{\pi} 3\)，\(c = \log_2 \sin \frac{2\pi}{5}\)，则（\quad）

\choices
    {\(a > b > c\)}
    {\(b > a > c\)}
    {\(c > a > b\)}
    {\(b > c > a\)}

\begin{answer}
A
\end{answer}

\begin{solution}
因为 \(1<3<\pi\)，且对数的底数 \(\pi>1\)，所以 \(0<b=\log_{\pi}3<1\). 又因为 \(a=\sqrt{2}>1\)，且 \(0<\sin\frac{2\pi}{5}<1\)，故 \(c=\log_2\sin\frac{2\pi}{5}<0\). 因此 \(a>b>c\)，故选 A.
\end{solution}
\end{problem}

\begin{problem}
“函数 \( f(x) \)（\(x \in \R\)）存在反函数”是“函数 \( f(x) \) 在 \(\R\) 上为增函数”的（\quad）

\choices
    {充分而不必要条件}
    {必要而不充分条件}
    {充分必要条件}
    {既不充分也不必要条件}

\begin{answer}
B
\end{answer}

\begin{solution}
函数在 \(\R\) 上为增函数时，对任意 \(x_1<x_2\) 都有 \(f(x_1)<f(x_2)\)，因此函数为单射，从而存在反函数. 反过来，存在反函数只说明函数为单射，不一定为增函数；例如 \(f(x)=-x\) 在 \(\R\) 上存在反函数，但它为减函数. 所以“存在反函数”是“在 \(\R\) 上为增函数”的必要而不充分条件，故选 B.
\end{solution}
\end{problem}

\begin{problem}
若点 \(P\) 到直线 \(x = -1\) 的距离比它到点 \((2,0)\) 的距离小 \(1\)，则点 \(P\) 的轨迹为（\quad）

\choices
    {圆}
    {椭圆}
    {双曲线}
    {抛物线}

\begin{answer}
D
\end{answer}

\begin{solution}
设 \(P=(x,y)\). 由题意，点 \(P\) 到点 \((2,0)\) 的距离比到直线 \(x=-1\) 的距离大 \(1\)，所以 \(\sqrt{(x-2)^2+y^2}=|x+1|+1\).

当 \(x\geqslant-1\) 时，右端为 \(x+2>0\)，平方得 \((x-2)^2+y^2=(x+2)^2\)，即 \(y^2=8x\). 该方程推出 \(x\geqslant0\)，与 \(x\geqslant-1\) 相容；反代可得两边都等于 \(x+2\)，故这些点均满足原方程.

当 \(x<-1\) 时，原方程化为 \(\sqrt{(x-2)^2+y^2}=-x\)，平方得 \(y^2=4x-4<0\)，没有实数解. 因而轨迹为抛物线 \(y^2=8x\)，故选 D.
\end{solution}
\end{problem}

\begin{problem}
若实数 \(x\)，\(y\) 满足 \(\begin{cases}
x - y + 1 \geqslant 0,\\
x + y \geqslant 0,\\
x \leqslant 0,
\end{cases}\) 则 \(z = 3^{x + 2y}\) 的最小值是（\quad）

\choices
    {\(0\)}
    {\(1\)}
    {\(\sqrt{3}\)}
    {\(9\)}

\begin{answer}
B
\end{answer}

\begin{solution}
由约束条件得 \(y\leqslant x+1\)，\(y\geqslant-x\)，且 \(x\leqslant0\). 可行的 \(x\) 还需满足 \(-x\leqslant x+1\)，所以 \(-\frac12\leqslant x\leqslant0\). 因为 \(3^t\) 在 \(\R\) 上递增，先求 \(x+2y\) 的最小值. 对固定的 \(x\)，取最小的 \(y=-x\)，于是 \(x+2y=-x\geqslant0\)，且在 \(x=0\) 时取得最小值 \(0\). 因此 \(z_{\min}=3^0=1\)，故选 B.
\end{solution}
\end{problem}

\begin{problem}
已知数列 \(\{a_{n}\}\) 对任意的 \(p\)，\(q \in \N^{*}\) 满足 \(a_{p + q} = a_p + a_q\)，且 \(a_2 = -6\)，那么 \(a_{10}\) 等于（\quad）

\choices
    {\(-165\)}
    {\(-33\)}
    {\(-30\)}
    {\(-21\)}

\begin{answer}
C
\end{answer}

\begin{solution}
取 \(p=q=1\)，得 \(a_2=2a_1\)，所以 \(a_1=-3\). 对任意正整数 \(n\)，取 \(p=n\)、\(q=1\)，有 \(a_{n+1}=a_n+a_1\). 用数学归纳法证明 \(a_n=-3n\)：当 \(n=1\) 时，\(a_1=-3\) 成立；若 \(a_n=-3n\)，则 \(a_{n+1}=a_n+a_1=-3n-3=-3(n+1)\). 因此 \(a_{10}=-3\times10=-30\)，故选 C.
\end{solution}
\end{problem}

\begin{problem}
过直线 \( y = x \) 上的一点作圆 \( (x - 5)^2 + (y - 1)^2 = 2 \) 的两条切线 \(l_1\)，\(l_2\)，当直线 \(l_1\)，\(l_2\) 关于 \( y = x \) 对称时，它们之间的夹角为（\quad）

\choices
    {\(30^{\circ}\)}
    {\(45^{\circ}\)}
    {\(60^{\circ}\)}
    {\(90^{\circ}\)}

\begin{answer}
C
\end{answer}

\begin{solution}
圆心为 \(O=(5,1)\)，半径为 \(r=\sqrt{2}\). 关于 \(y=x\) 的反射把 \(O\) 变为 \(O'=(1,5)\). 由于两条切线关于 \(y=x\) 对称，且交点在 \(y=x\) 上，反射后每条切线仍与原圆相切，因此它们是以 \(O\)、\(O'\) 为圆心、半径均为 \(\sqrt2\) 的两个等圆的两条内公切线. 两内公切线的交点是两圆心的中点 \(M=(3,3)\).

有 \(MO=\sqrt{(5-3)^2+(1-3)^2}=2\sqrt2\). 设从 \(M\) 引出的其中一条切线与直线 \(MO\) 的锐角为 \(\alpha\)，则在切点处半径垂直切线，所以 \(\sin\alpha=\frac{\sqrt2}{2\sqrt2}=\frac12\)，从而 \(\alpha=30^{\circ}\). 两条切线关于 \(MO\) 对称，故它们的较小夹角为 \(2\alpha=60^{\circ}\)，故选 C.
\end{solution}
\end{problem}

\begin{problem}
如图，动点 \(P\) 在正方体 \(ABCD - A_{1}B_{1}C_{1}D_{1}\) 的对角线 \(BD_{1}\) 上. 过点 \(P\) 作垂直于平面 \(BB_{1}D_{1}D\) 的直线. 与正方体表面相交于 \(MN\). 设 \(BP = x\)，\(MN = y\)，则函数 \(y = f(x)\) 的图象大致是（\quad）

\bitmapfigure[width=0.84\linewidth]{img/2008/beijing_science/q8_fig1.png}

\choices
    {\choicebitmap[width=0.15\paperwidth]{img/2008/beijing_science/q8_optA.png}}
    {\choicebitmap[width=0.15\paperwidth]{img/2008/beijing_science/q8_optB.png}}
    {\choicebitmap[width=0.15\paperwidth]{img/2008/beijing_science/q8_optC.png}}
    {\choicebitmap[width=0.15\paperwidth]{img/2008/beijing_science/q8_optD.png}}

\begin{answer}
B
\end{answer}

\begin{solution}
设正方体棱长为 \(a\)，令 \(t=\dfrac{BP}{BD_1}\)，则 \(0\leqslant t\leqslant1\)，且 \(BP=a\sqrt3t\). 取坐标 \(A=(0,0,0)\)，\(B=(a,0,0)\)，\(C=(a,a,0)\)，\(D=(0,a,0)\)，\(A_1=(0,0,a)\)，\(B_1=(a,0,a)\)，\(D_1=(0,a,a)\). 于是 \(P=(a(1-t),at,at)\).

平面 \(BB_1D_1D\) 的方程为 \(x+y=a\)，其法向量可取 \((1,1,0)\). 过 \(P\) 的垂线保持 \(z=at\)，在正方形截面 \(0\leqslant x\)，\(y\leqslant a\) 中沿方向 \((1,1)\) 通过 \(P\). 当 \(0\leqslant t\leqslant\frac12\) 时，直线段 \(MN\) 的两个端点为 \((a(1-2t),0,at)\) 和 \((a,2at,at)\)，故 \(MN=2\sqrt2at\)；当 \(\frac12\leqslant t\leqslant1\) 时，两个端点为 \((0,a(2t-1),at)\) 和 \((a(2-2t),a,at)\)，故 \(MN=2\sqrt2a(1-t)\). 因而 \(MN=2\sqrt2a\min\{t,1-t\}\).

代入 \(t=\dfrac{x}{a\sqrt3}\)，得到 \(y=2\sqrt2a\min\left\{\dfrac{x}{a\sqrt3},1-\dfrac{x}{a\sqrt3}\right\}\). 因而 \(y\) 从 \(x=0\) 时的 \(0\) 线性增加，在 \(x=\frac{a\sqrt3}{2}\) 处达到最大值，再线性减小到 \(x=a\sqrt3\) 时的 \(0\). 图象是关于 \(x=\frac{a\sqrt3}{2}\) 对称的折线三角形，故选 B.
\end{solution}
\end{problem}

\section{填空题}

\begin{problem}
已知 \((a - \i)^2 = 2\i\)，其中 \(\i\) 是虚数单位，那么实数 \(a = \)\fillinblank{}.

\begin{answer}
\(-1\)
\end{answer}

\begin{solution}
展开左边得 \((a-\i)^2=a^2-1-2a\i\). 与 \(2\i\) 比较实部和虚部，得到 \(a^2-1=0\) 以及 \(-2a=2\). 后一个等式给出 \(a=-1\)，且满足前一个等式，故填 \(-1\).
\end{solution}
\end{problem}

\begin{problem}
已知向量 \(\bs{a}\) 与 \(\bs{b}\) 的夹角为 \(120^{\circ}\)，且 \(|\bs{a}| = |\bs{b}| = 4\)，那么 \(\bs{b} \cdot (2\bs{a} + \bs{b})\) 的值为\fillinblank{}.

\begin{answer}
\(0\)
\end{answer}

\begin{solution}
由向量夹角公式，\(\bs{a}\cdot\bs{b}=|\bs{a}||\bs{b}|\cos120^{\circ}=4\cdot4\cdot\left(-\frac12\right)=-8\). 因此 \(\bs{b}\cdot(2\bs{a}+\bs{b})=2\bs{a}\cdot\bs{b}+\bs{b}\cdot\bs{b}=2(-8)+4^2=0\)，故填 \(0\).
\end{solution}
\end{problem}

\begin{problem}
若 \(\left(x^{2} + \frac{1}{x^{3}}\right)^{n}\) 展开式的各项系数之和为 \(32\)，则 \(n =\fillinblank{}\)，其展开式中的常数项为\fillinblank{}. （用数字作答）

\begin{answer}
\(5\)，\(10\)
\end{answer}

\begin{solution}
令展开式中的 \(x^2\) 取 \(k\) 次，则通项为 \(\mathrm C_n^k(x^2)^k\left(\frac1{x^3}\right)^{n-k}=\mathrm C_n^k x^{5k-3n}\). 各项系数之和等于令 \(x=1\) 时的值，所以 \(2^n=32\)，得到 \(n=5\). 常数项要求 \(5k-3n=0\)，代入 \(n=5\) 得 \(k=3\)，其系数为 \(\mathrm C_5^3=10\). 故填 \(5\)，\(10\).
\end{solution}
\end{problem}

\begin{problem}
如图，函数 \( f(x) \) 的图象是折线段 \(ABC\)，其中 \(A\)，\(B\)，\(C\) 的坐标分别为 \((0,4)\)，\((2,0)\)，\((6,4)\)，则 \( f(f(0)) = \)\fillinblank{}；\(\displaystyle\lim_{\Delta x \to 0} \frac{f(1 + \Delta x) - f(1)}{\Delta x} = \)\fillinblank{}. （用数字作答）

\bitmapfigure[width=6.0cm]{img/2008/beijing_science/q12_fig1.png}

\begin{answer}
\(2\)，\(-2\)
\end{answer}

\begin{solution}
在 \(A(0,4)\)、\(B(2,0)\) 所在的线段上，斜率为 \(-2\)，所以 \(f(x)=4-2x\)（\(0\leqslant x\leqslant2\)）；在 \(B(2,0)\)、\(C(6,4)\) 所在的线段上，斜率为 \(1\)，所以 \(f(x)=x-2\)（\(2\leqslant x\leqslant6\)）. 于是 \(f(0)=4\)，\(f(f(0))=f(4)=2\).

当 \(\Delta x\) 足够小时，\(1+\Delta x\in[0,2]\)，故 \(\frac{f(1+\Delta x)-f(1)}{\Delta x}=\frac{[4-2(1+\Delta x)]-[4-2]}{\Delta x}=-2\). 所以极限为 \(-2\)，故填 \(2\)，\(-2\).
\end{solution}
\end{problem}

\begin{problem}
已知函数 \( f(x) = x^2 - \cos x \)，对于 \( \left[-\frac{\pi}{2}, \frac{\pi}{2}\right] \) 上的任意 \(x_1\)，\(x_2\)，有如下条件：

\begin{circlelist}
\item \( x_1 > x_2 \)；
\item \( x_1^2 > x_2^2 \)；
\item \( |x_1| > x_2 \).
\end{circlelist}

其中能使 \( f(x_1) > f(x_2) \) 恒成立的条件序号是\fillinblank{}.

\begin{answer}
\circled{2}
\end{answer}

\begin{solution}
令 \(g(t)=t^2-\cos t\)（\(0\leqslant t\leqslant\frac\pi2\)），则 \(g'(0)=0\)，且 \(g'(t)=2t+\sin t>0\)（\(0<t\leqslant\frac\pi2\)）. 因而对任意 \(0\leqslant s<t\leqslant\frac\pi2\)，有 \(g(t)-g(s)=\int_s^t(2u+\sin u)\,\mathrm du>0\)，所以 \(g\) 在该区间上严格递增；又 \(f(x)=g(|x|)\).

条件\(\circled{2}\)给出 \(|x_1|>|x_2|\)，从而 \(f(x_1)=g(|x_1|)>g(|x_2|)=f(x_2)\)，因此条件\(\circled{2}\)充分. 条件\(\circled{1}\)不充分，例如 \(x_1=0\)、\(x_2=-\frac12\) 时 \(x_1>x_2\)，但 \(f(0)=-1\)，而 \(f\left(-\frac12\right)=\frac14-\cos\frac12>-\frac34>-1\). 条件\(\circled{3}\)也不充分，例如 \(x_1=-1\)、\(x_2=-\frac\pi2\) 时 \(|x_1|>x_2\)，但 \(f(-1)=1-\cos1<1<\frac{\pi^2}{4}=f\left(-\frac\pi2\right)\). 因此填 \(2\).
\end{solution}
\end{problem}

\begin{problem}
某校数学课外小组在坐标纸上，为学校的一块空地设计植树方案如下：第 \(k\) 棵树种植在点 \(P_{k}(x_{k},y_{k})\) 处，其中 \(x_{1} = 1\)，\(y_{1} = 1\)，当 \(k \geqslant 2\) 时，\(\begin{cases}
x_{k} = x_{k - 1} + 1 - 5\left[ T\left(\frac{k - 1}{5}\right) - T\left(\frac{k - 2}{5}\right) \right],\\
y_{k} = y_{k - 1} + T\left(\frac{k - 1}{5}\right) - T\left(\frac{k - 2}{5}\right).
\end{cases}\) \(T(a)\) 表示非负实数 \(a\) 的整数部分，例如 \(T(2.6) = 2\)，\(T(0.2) = 0\). 按此方案，第 \(6\) 棵树种植点的坐标应为\fillinblank{}；第 \(2008\) 棵树种植点的坐标应为\fillinblank{}.

\begin{answer}
\((1,2)\)，\((3,402)\)
\end{answer}

\begin{solution}
记 \(d_k=T\left(\frac{k-1}{5}\right)-T\left(\frac{k-2}{5}\right)\). 当 \(k\geqslant2\) 时，\(d_k=1\) 当且仅当 \(k=5j+1\)（\(j\) 为正整数），其余情形 \(d_k=0\). 因而通常每次使 \(x\) 增加 \(1\)，使 \(y\) 不变；在 \(k=5j+1\) 时，\(x\) 的增量为 \(-4\)，\(y\) 的增量为 \(1\).

从 \(P_1=(1,1)\) 到 \(P_6\)，只有 \(k=6\) 时 \(d_k=1\)，故 \(P_6=(1,2)\). 在 \(k=2\)，\(\ldots\)，\(2008\) 中，满足 \(k=5j+1\) 的 \(k\) 有 \(\left\lfloor\frac{2007}{5}\right\rfloor=401\) 个. 若没有这些特殊增量，\(x_{2008}=1+2007=2008\)；每次特殊增量相对于通常的增量 \(1\) 少 \(5\)，所以 \(x_{2008}=2008-5\cdot401=3\)，而 \(y_{2008}=1+401=402\). 故填 \((1,2)\)，\((3,402)\).
\end{solution}
\end{problem}

\section{解答题}

\begin{problem}
已知函数 \( f(x) = \sin^2 \omega x + \sqrt{3}\sin \omega x\sin \left(\omega x + \frac{\pi}{2}\right) \)（\(\omega > 0\)）的最小正周期为 \( \pi \).

\begin{enumerate}
\item 求 \(\omega\) 的值；
\item 求函数 \( f(x) \) 在区间 \( \left[0, \frac{2\pi}{3}\right] \) 上的取值范围.
\end{enumerate}

\begin{solution}
\begin{enumerate}
\item 令 \(u=\omega x\). 因为 \(\sin\left(u+\frac\pi2\right)=\cos u\)，所以
\(\displaystyle f(x)=\sin^2u+\sqrt3\sin u\cos u=\frac12+\frac{\sqrt3}{2}\sin2u-\frac12\cos2u=\frac12+\sin\left(2u-\frac\pi6\right)\).
因此 \(f(x)\) 的最小正周期为 \(\frac\pi\omega\). 由题设 \(\frac\pi\omega=\pi\)，且 \(\omega>0\)，故 \(\omega=1\).
\item 此时 \(f(x)=\frac12+\sin\left(2x-\frac\pi6\right)\). 当 \(0\leqslant x\leqslant\frac{2\pi}{3}\) 时，令 \(t=2x-\frac\pi6\)，则 \(-\frac\pi6\leqslant t\leqslant\frac{7\pi}{6}\). 该区间包含 \(t=\frac\pi2\)，故函数最大值为 \(\frac12+1=\frac32\). 区间两端的正弦值均为 \(-\frac12\)，而区间内不含正弦值为 \(-1\) 的点 \(-\frac\pi2+2k\pi\)，故最小值为 \(\frac12-\frac12=0\). 取值范围为 \(\left[0,\frac32\right]\).
\end{enumerate}
\end{solution}
\end{problem}

\begin{problem}
如图，在三棱锥 \(P - ABC\) 中，\(AC = BC = 2\)，\(\angle ACB = 90^{\circ}\)，\(AP = BP = AB\)，\(PC \perp AC\).

\begin{enumerate}
\item 求证：\(PC \perp AB\)；
\item 求二面角 \(B - AP - C\) 的大小；
\item 求点 \(C\) 到平面 \(APB\) 的距离.
\end{enumerate}

\bitmapfigure[width=6.0cm]{img/2008/beijing_science/q16_fig1.png}

\begin{solution}
\begin{enumerate}
\item 因为 \(AC=BC\)，点 \(C\) 在 \(AB\) 的垂直平分面上；又因为 \(AP=BP\)，点 \(P\) 也在 \(AB\) 的垂直平分面上. 因而直线 \(PC\) 在这个平面内，而该平面垂直于 \(AB\)，所以 \(PC\perp AB\).
\item 取 \(C=(0,0,0)\)，\(A=(2,0,0)\)，\(B=(0,2,0)\). 由 \(PC\perp AC\)，可设 \(P=(0,y,z)\). 由 \(AP=BP\) 得 \(4+y^2+z^2=(y-2)^2+z^2\)，从而 \(y=0\). 又 \(AP=AB=2\sqrt2\)，所以 \(4+z^2=8\)，即 \(z^2=4\). 按图取 \(P=(0,0,2)\).

平面 \(CAP\) 的方程为 \(y=0\)，平面 \(ABP\) 的方程为 \(x+y+z=2\). 两平面的法向量可分别取 \(\bs{n}_1=(0,1,0)\)、\(\bs{n}_2=(1,1,1)\)，故二面角的锐角 \(\theta\) 满足 \(\cos\theta=\frac{|\bs{n}_1\cdot\bs{n}_2|}{|\bs{n}_1||\bs{n}_2|}=\frac1{\sqrt3}\). 因此 \(\sin\theta=\sqrt{\frac23}=\frac{\sqrt6}{3}\)，二面角大小为 \(\arcsin\frac{\sqrt6}{3}\).
\item 由平面 \(ABP:x+y+z=2\) 的点到平面距离公式，点 \(C=(0,0,0)\) 到该平面的距离为 \(\frac{|0+0+0-2|}{\sqrt{1^2+1^2+1^2}}=\frac2{\sqrt3}=\frac{2\sqrt3}{3}\).
\end{enumerate}
\end{solution}
\end{problem}

\begin{problem}
甲、乙等五名奥运志愿者被随机地分到 A，B，C，D 四个不同的岗位服务，每个岗位至少有一名志愿者.

\begin{enumerate}
\item 求甲、乙两人同时参加 A 岗位服务的概率；
\item 求甲、乙两人不在同一个岗位服务的概率；
\item 设随机变量 \( \xi \) 为这五名志愿者中参加 A 岗位服务的人数，求 \( \xi \) 的分布列.
\end{enumerate}

\begin{solution}
\begin{enumerate}
\item 因为 \(5\) 名志愿者分到 \(4\) 个岗位且每个岗位至少一人，所以人数分配必为 \(2\)，\(1\)，\(1\)，\(1\). 先选出二人所在的岗位，再选出这二人，最后把其余三人分别安排到其余三个岗位，等可能的分配总数为 \(4\mathrm C_5^2 3!=4\times10\times6=240\). 甲、乙同时在 A 岗位时，二人必须正好是 A 岗位的两人，其余三人分别到 B，C，D，有 \(3!\) 种，所以概率为 \(\frac{3!}{240}=\frac1{40}\).
\item 甲、乙在同一岗位时，他们必为唯一的二人岗位. 先选这个岗位，再把其余三人分别安排到另外三个岗位，共有 \(4\cdot3!=24\) 种，因此 \(P(\text{甲、乙同岗})=\frac{24}{240}=\frac1{10}\)，从而 \(P(\text{甲、乙不同岗})=1-\frac1{10}=\frac9{10}\).
\item A 岗位人数只能取 \(1\) 或 \(2\). 当 \(\xi=2\) 时，先从五人中选出到 A 岗位的二人，再把其余三人分别安排到 B，C，D，有 \(\mathrm C_5^2 3!=60\) 种，故 \(P(\xi=2)=\frac{60}{240}=\frac14\). 因此 \(P(\xi=1)=1-\frac14=\frac34\). 分布列为
\begin{center}
\begin{tabular}{c|cc}
\(\xi\)&\(1\)&\(2\)\\ \hline
\(P\)&\(\frac34\)&\(\frac14\)
\end{tabular}
\end{center}
\end{enumerate}
\end{solution}
\end{problem}

\begin{problem}
已知函数 \( f(x) = \frac{2x - b}{(x - 1)^2} \)，求导函数 \( f'(x) \)，并确定 \( f(x) \) 的单调区间.

\begin{solution}
\(f'(x)=\frac{2(x-1)^2-2(2x-b)(x-1)}{(x-1)^4}=-\frac{2\bigl[x-(b-1)\bigr]}{(x-1)^3}\)，\(x\ne1\).
令 \(c=b-1\). 当 \(b<2\)（即 \(c<1\)）时，\(f'(x)<0\) 在 \((-\infty,c)\) 与 \((1,+\infty)\) 上成立，\(f'(x)>0\) 在 \((c,1)\) 上成立，所以 \(f(x)\) 在 \((-\infty,b-1)\)、\((1,+\infty)\) 上单调递减，在 \((b-1,1)\) 上单调递增. 当 \(b>2\)（即 \(c>1\)）时，\(f'(x)<0\) 在 \((-\infty,1)\) 与 \((c,+\infty)\) 上成立，\(f'(x)>0\) 在 \((1,c)\) 上成立，所以 \(f(x)\) 在 \((-\infty,1)\)、\((b-1,+\infty)\) 上单调递减，在 \((1,b-1)\) 上单调递增. 当 \(b=2\) 时，\(f(x)=\frac2{x-1}\)，且 \(f'(x)=-\frac2{(x-1)^2}<0\)，所以在 \((-\infty,1)\)、\((1,+\infty)\) 上均单调递减.
\end{solution}
\end{problem}

\begin{problem}
已知菱形 \(ABCD\) 的顶点 \(A\)，\(C\) 在椭圆 \(x^{2} + 3y^{2} = 4\) 上，对角线 \(BD\) 所在直线的斜率为 \(1\).

\begin{enumerate}
\item 当直线 \(BD\) 过点 \((0,1)\) 时，求直线 \(AC\) 的方程；
\item 当 \(\angle ABC = 60^{\circ}\) 时，求菱形 \(ABCD\) 面积的最大值.
\end{enumerate}

\begin{solution}
\begin{enumerate}
\item 设 \(O\) 为菱形两条对角线的交点. 由 \(BD\) 的斜率为 \(1\)，\(AC\perp BD\)，所以可设 \(O=(u,v)\)，且 \(v=u+1\). 取非零实数 \(t\)，令 \(A=O+t(1,-1)\)、\(C=O-t(1,-1)\). 因为 \(A\)，\(C\) 都在椭圆 \(x^2+3y^2=4\) 上，比较两式 \(u+t\)，\(v-t\) 与 \(u-t\)，\(v+t\) 代入椭圆方程所得结果，得 \(4t(u-3v)=0\)，所以 \(u=3v\). 联立 \(v=u+1\)，得 \(u=-\frac32\)、\(v=-\frac12\). 因而 \(AC\) 的方向为 \((1,-1)\)，且 \(x+y=u+v=-2\)，故直线 \(AC\) 为 \(x+y+2=0\).
\item 设菱形边长为 \(s\). 由 \(AB=BC=s\) 及 \(\angle ABC=60^{\circ}\)，在三角形 \(ABC\) 中由余弦定理得 \(AC^2=s^2+s^2-2s^2\cos60^{\circ}=s^2\). 又因为菱形对角线互相垂直且满足 \(AC^2+BD^2=2(AB^2+BC^2)\)，所以 \(BD^2=3s^2\). 菱形面积为 \(\frac12AC\cdot BD=\frac{\sqrt3}{2}AC^2\).

设斜率为 \(-1\) 的直线 \(AC\) 写为 \(x+y=c\)，令 \(h=\frac c2\)，其上点可表示为 \((h+r,h-r)\). 代入椭圆方程得 \(4r^2-4hr+4h^2-4=0\). 设两个交点对应的参数为 \(r_1\)，\(r_2\)，由求根公式或根的判别式，\[
r_1-r_2=\sqrt{(r_1+r_2)^2-4r_1r_2}=\sqrt{4-3h^2}.
\]
由于参数方向向量 \((1,-1)\) 的长度为 \(\sqrt2\)，所以弦长
\(AC=\sqrt2\sqrt{4-3h^2}\leqslant2\sqrt2\)，等号在 \(h=0\)，即 \(AC\) 为直线 \(x+y=0\) 与椭圆的弦时取得. 当 \(h=0\) 时，椭圆与直线 \(x+y=0\) 的交点可取 \(A=(1,-1)\)、\(C=(-1,1)\)，且 \(AC=2\sqrt2\). 取 \(B=(\sqrt3,\sqrt3)\)、\(D=(-\sqrt3,-\sqrt3)\)，则 \(BD=2\sqrt6=\sqrt3\,AC\)，并且 \(AB^2=BC^2=AC^2=8\)，由关于原点的中心对称性还有 \(CD=DA=AB\)，所以 \(ABCD\) 确实为菱形；由余弦定理，\(\cos\angle ABC=\frac{AB^2+BC^2-AC^2}{2AB\cdot BC}=\frac12\)，故 \(\angle ABC=60^{\circ}\). 因而菱形面积最大值为 \(\frac{\sqrt3}{2}(2\sqrt2)^2=4\sqrt3\).
\end{enumerate}
\end{solution}
\end{problem}

\begin{problem}
对于每项均是正整数的数列 \(A\)： \(a_1\)，\(a_2\)，\(\cdots\)，\(a_n\)，定义变换 \(T_1\) 将数列 \(A\) 变换成数列 \(T_1\!\left(A\right)\)： \(n\)，\(a_1 - 1\)，\(a_2 - 1\)，\(\cdots\)，\(a_n - 1\). 对于每项均是非负整数的数列 \(B\)： \(b_1\)，\(b_2\)，\(\cdots\)，\(b_m\)，定义变换 \(T_2\) 将数列 \(B\) 各项从大到小排列，然后去掉所有为零的项，得到数列 \(T_2\!\left(B\right)\)；又定义 \(S\!\left(B\right) = 2(b_1 + 2b_2 + \cdots + mb_m) + b_1^2 + b_2^2 + \cdots + b_m^2\). 设 \(A_0\) 是每项均为正整数的有穷数列，令 \(A_{k+1} = T_2\!\left(T_1\!\left(A_k\right)\right)\)（\(k = 0\)，\(1\)，\(2\)，\(\cdots\)）.

\begin{enumerate}
\item 如果数列 \(A_0\) 为 \(5\)，\(3\)，\(2\)，写出数列 \(A_1\)，\(A_2\)；
\item 对于每项均是正整数的有穷数列 \(A\)，证明 \(S\!\left(T_1\!\left(A\right)\right) = S\!\left(A\right)\)；
\item 证明：对于任意给定的每项均为正整数的有穷数列 \(A_0\)，存在正整数 \(K\)，当 \(k \geqslant K\) 时，\(S(A_{k+1}) = S(A_k)\).
\end{enumerate}

\begin{solution}
\begin{enumerate}
\item \(T_1(A_0)=(3,4,2,1)\)，按从大到小排列并去掉零项后，\(A_1=(4,3,2,1)\). 再有 \(T_1(A_1)=(4,3,2,1,0)\)，所以 \(A_2=(4,3,2,1)\).
\item 设 \(A=(a_1,a_2,\ldots,a_n)\)，则 \(T_1(A)=(n,a_1-1,a_2-1,\ldots,a_n-1)\). 直接展开 \(S\) 得
\(S\left(T_1(A)\right)=2\left[n+\sum_{i=1}^{n}(i+1)(a_i-1)\right]+n^2+\sum_{i=1}^{n}(a_i-1)^2\).
展开后，\(a_i\) 的一次项为 \(2(i+1)a_i-2a_i=2ia_i\)，平方项为 \(\sum_{i=1}^n a_i^2\)，常数项为 \(2n-2\sum_{i=1}^{n}(i+1)+n^2+n\).
又 \(\sum_{i=1}^{n}(i+1)=\frac{n(n+3)}2\)，所以该常数项等于 \(2n-n(n+3)+n^2+n=0\). 因此 \(S\left(T_1(A)\right)=2\sum_{i=1}^{n}ia_i+\sum_{i=1}^{n}a_i^2=S(A)\).
\item 对任意非负整数数列 \(B=(b_1,\ldots,b_m)\)，交换相邻的逆序项 \(b_i<b_{i+1}\) 后，平方和不变，而加权和的变化为
\[
i b_i+(i+1)b_{i+1}-\bigl(i b_{i+1}+(i+1)b_i\bigr)=b_{i+1}-b_i>0
\]
这里左边是交换前的加权和，右边是交换后的加权和，所以每次交换都会使 \(S\) 严格减小；将所有逆序项交换后，数列按从大到小排列，故排序不会增加 \(S\). 排序后所有零项位于末尾，去掉零项也不改变 \(S\)，从而 \(S(T_2(B))\leqslant S(B)\). 结合第（2）问，对 \(A_k\) 有
\[
S(A_{k+1})=S\left(T_2(T_1(A_k))\right)\leqslant S\left(T_1(A_k)\right)=S(A_k)
\]
由于每个 \(A_k\) 都是非负整数项数列去掉零项后得到的正整数项数列，故 \(S(A_k)\) 是非负整数. 单调不增的非负整数列不可能无限次严格下降，因此从某个正整数 \(K\) 起不再下降，即对所有 \(k\geqslant K\)，都有 \(S(A_{k+1})=S(A_k)\).
\end{enumerate}
\end{solution}
\end{problem}
